% !TEX program = xelatex
% The research corpus (ci/research.json): a XeLaTeX paper with fontspec
% and unicode-math on TeX Live's OpenType fonts, by name and by file.
\documentclass[11pt,a4paper]{article}
\usepackage{fontspec}
\usepackage{unicode-math}
\setmainfont{TeX Gyre Termes}
\setsansfont{TeX Gyre Heros}[Scale=MatchLowercase]
\setmonofont{TeX Gyre Cursor}[Scale=MatchLowercase]
\setmathfont{TeX Gyre Termes Math}
\newfontfamily\libertinus{LibertinusSerif}[Extension=.otf,UprightFont=*-Regular,ItalicFont=*-Italic,BoldFont=*-Bold]
\newfontfamily\pagella{TeX Gyre Pagella}[Numbers=OldStyle]
\usepackage{booktabs}
\usepackage[hidelinks]{hyperref}
\title{Unicode Text and Mathematics in XeLaTeX}
\author{J. Typographer}
\date{2026}
\begin{document}
\maketitle
\section{Text}
Unicode input — “quotes”, ‘single’, en–dash, em—dash, ellipsis…, ligatures
fi fl ffi, accents àéîõü ç ñ ø å, and \textsf{sans serif}, \texttt{mono},
\textit{italic}, \textbf{bold}, \textsc{Small Caps}.

{\libertinus Libertinus by its file name: \textit{italic} and \textbf{bold}.}

{\pagella Old-style figures in Pagella: 0123456789.}
\section{Mathematics}
With unicode-math, $α + β = γ$, $∀x ∈ ℝ: x² ≥ 0$ and
\[
  \int_{-\infty}^{\infty} e^{-x^2}\,dx = \sqrt{\pi}, \qquad
  \mathbf{A}\mathbf{x} = \mathbf{b}, \qquad
  \mathcal{F}\{f\}(\xi) = \int f(x)\,e^{-2\pi i x \xi}\,dx .
\]
\begin{equation}
  \sum_{k=0}^{n} \binom{n}{k} = 2^n , \qquad \mathbb{E}[X] = \sum_x x\,P(X=x).
\end{equation}
\section{Table}
\begin{tabular}{lrr}
\toprule
Font & Glyphs & Scripts \\
\midrule
TeX Gyre Termes & 1200 & Latin \\
Libertinus Serif & 2700 & Latin, Greek, Cyrillic \\
\bottomrule
\end{tabular}
\end{document}
